% 重构稿：硬可行域约束下的有限时域CMDP与学习型端边云调度。
\section{TIDE-Placement：端边云自适应推理调度}
\label{sec:tide_placement}

\subsection{设计目标与约束决策架构}

端边云放置需要处理长程任务中的跨步骤资源取舍。将某个复杂节点放到高能力资源上，可能提高当前结果质量，却消耗后续节点所需的时间与费用预算；在本地快速响应，也可能因结果不足而触发更多修订。因此，资源选择不仅依赖当前任务复杂度，还应考虑剩余任务、可用预算、服务负载以及既往尝试的执行反馈。

本文将TIDE-Placement组织为\textbf{硬可行域约束下的学习型资源调度架构}。可信域、权限、模型能力和当前可用性先确定允许动作，策略再在允许范围内优化任务质量与累计资源消耗。其建模借鉴约束强化学习将收益与约束预算分开定义的思想\cite{tide_cpo}；具体采用有限时域约束决策过程与拉格朗日Actor--Critic设计，不将其等同于CPO算法，也不继承后者的理论保证。

架构包括状态编码器、可行域过滤器、策略与价值评估器、约束预算控制器、执行反馈器五个部分。BGE类编码器与MLP可分别承载任务语义表示和策略网络；规则调度器作为冷启动及学习策略不可用时的回退。只有依据奖励和约束轨迹更新策略，才能称为强化学习，已有监督分类路由不能仅通过改写公式获得这一性质。

\begin{figure}[htbp]
\centering
\fbox{\parbox{0.92\textwidth}{\centering
任务及剩余预算 $+$ 资源监控 $\rightarrow$ 状态编码\\[3pt]
$\downarrow$\qquad\qquad\qquad\qquad $\downarrow$\\[3pt]
硬可行域过滤 $\rightarrow$ 掩码策略与价值评估\\[3pt]
$\rightarrow$ 阶段计划与安全处理 $\rightarrow$ 执行器分发\\[3pt]
$\rightarrow$ 质量、时延、费用与失败反馈\\[3pt]
$\longrightarrow$ 更新状态与预算 $\longrightarrow$ 轨迹及策略训练}}
\caption{硬约束与累计预算共同作用的学习型调度架构}
\label{fig:tide-placement-loop}
\end{figure}

\subsection{任务、资源与调度动作的统一表示}

\paragraph{状态构造}
以一次可调度的推理请求或失败后的重新申请为决策步$t$，任务结束、不可恢复失败或预算终止为回合结束。有限时域上限为$H$，包含重试次数，避免通过无限等待规避任务失败。状态表示为
\begin{equation}
s_t=\langle z_t,p_t,\omega_t,b_t,h_t,v_t\rangle,
\label{eq:placement_state}
\end{equation}
其中，$z_t$为任务语义及当前输入的表示，$p_t$为任务画像，$\omega_t$为资源状态，$b_t$为费用、时间和尝试等剩余预算，$h_t$为必要历史统计，$v_t$为计划及约束版本。使用共享状态的最小必要投影，不将整个长期对话输入调度器。

任务画像包含实时性、敏感度、复杂度、任务类型、预期产物质量和输入长度；显式元数据优先于模型推断，缺失字段由门控器补充。资源状态记录端点健康度、可用模型、上下文容量、负载、队列、近期时延、限流和错误统计。可信属性来自部署和授权配置，不能把“云”天然视为不可信，也不能把“边缘”天然视为可信。

资源画像应针对具体端点和模型配置。端、边、云只是位置标签：设备可能更快但能力有限，边缘可能适合内部数据处理，云侧可能提供更高能力模型，实际比较以部署参数和探测结果为准。参考时延与成本权重是配置值，不能写成实测硬件收益。

状态未必完整描述环境。网络和服务负载可能不可观测，式~\eqref{eq:placement_state}是用于决策的近似充分统计；必要时加入滑动窗口或历史编码。状态表示是否足以支持稳定决策，需要通过分布变化和负载扰动实验验证。

\paragraph{动作与阶段计划}
资源动作定义为
\begin{equation}
a_t=\langle r_t,m_t,\xi_t\rangle\in\mathcal{A}(s_t),
\end{equation}
其中$r_t$是具体资源端点，$m_t$为可用模型配置，$\xi_t$为经批准的推理参数。动作空间规模需受到限制，以便估计质量、时延与费用。无可行推理资源时，控制器返回等待、拆分请求或需要复核的明确状态；这些处理受时限和次数约束，不作为任意长的无成本动作。

调度动作进一步生成任务级阶段计划$\Pi_t$：任务画像构建、必要的可信预处理、一次主推理及可选产物检查。典型路径为
\begin{equation}
\text{本地画像}\rightarrow\text{可信预处理（按需）}
\rightarrow\text{选定端点主推理}\rightarrow\text{质量检查}.
\end{equation}
本节讨论的是任务阶段放置，不等同于将单个神经网络的层或张量切分到多台设备。阶段名称和资源位置显式记录，使计划可以被执行器和审计模块共同读取。

\subsection{硬可行域与有限时域约束建模}

\paragraph{不能被收益补偿的硬约束}
可行域过滤器先检查动作是否合法：
\begin{equation}
\mathcal{A}_{\mathrm{safe}}(s_t)=
\left\{a\mid
\begin{array}{l}
\operatorname{Available}(a,s_t)=1,\\
\operatorname{Authorized}(a,s_t)=1,\\
\operatorname{TrustOK}(a,s_t)=1,\\
\operatorname{CapabilityOK}(a,s_t)=1
\end{array}\right\}.
\label{eq:placement_safe_set}
\end{equation}
能力条件包括上下文容量、必要模型接口和阶段输出要求。对必须达到的能力门槛严格过滤；可选的质量偏好留给策略优化。将高复杂任务送到能力不足的本地模型必须经过显式任务降级或拆分，不能作为满足原始要求的成功回退。

敏感载荷仅能流向被当前策略授权的处理域。若采用脱敏以允许外部处理，必须验证输入、消息、附件及请求元数据的实际发送载荷均符合相应策略，不能仅因存在脱敏阶段就自动允许外传。安全预处理后和请求发送前再次核验，以处理策略、资源和载荷在调度后发生变化的情况。

\begin{figure}[htbp]
\centering
\fbox{\parbox{0.92\textwidth}{\centering
敏感性与授权策略 $+$ 具体端点可信属性\\[3pt]
$\rightarrow$ 生成允许动作 $\rightarrow$ 策略选取候选\\[3pt]
$\rightarrow$ 可信预处理与实际载荷核验\\[3pt]
$\rightarrow$ 通过后执行；不满足则阻断、改选或复核}}
\caption{独立于学习收益的可信域与发送前校验}
\label{fig:placement-trusted-path}
\end{figure}

\paragraph{收益与约束成本}
有限时域约束马尔可夫决策过程用状态、动作、未知转移、收益和约束成本描述调度。收益$r_t$由可核验任务进展及阶段质量构成，终止时加入最终交付评价。重复完成同一产物不能重复获得进展奖励；提前放弃按未完成任务评价，不能仅因省费用而被评为优质策略。模型评审可补充语义质量，但需结合规则和真实证据校准。

定义费用成本$c_{\mathrm{fee},t}$、经过时间$c_{\mathrm{time},t}$及失败尝试$c_{\mathrm{fail},t}$。费用包含实际调用及安全处理、重试等可计费开销；时间按墙钟区间的增量累计，并行运行不能简单累加各节点耗时作为端到端时延。能耗仅在有测量或经过校准的估计器时加入，不能以费用权重冒充功耗。

策略目标为
\begin{equation}
\begin{aligned}
\max_{\pi}\quad &J_R(\pi)=
\mathbb{E}_{\pi}\!\left[\sum_{t=0}^{H-1}r_t\right]\\
\mathrm{s.t.}\quad &J_i(\pi)=
\mathbb{E}_{\pi}\!\left[\sum_{t=0}^{H-1}c_{i,t}\right]\leq d_i,
\quad i\in\{\mathrm{fee,time,fail}\},\\
&\pi(a\mid s)=0,\quad a\notin\mathcal{A}_{\mathrm{safe}}(s).
\end{aligned}
\label{eq:placement_cmdp}
\end{equation}
$d_i$为对应任务类别的预算。若不同任务预算不同，应将预算纳入状态并分别报告可行性，不用混合平均值掩盖困难任务超限。训练可按各任务预算归一化成本，最终报告仍使用实际费用、时间和次数。

这里的累计约束是\textbf{期望约束}，不能保证每条任务轨迹都不超限。运行时另设剩余预算检查：若可预估开销将超过剩余额度则阻断或拆分，并对单次调用设置费用或超时上限。网络时延和失败具有随机性，预估不能提供确定性截止时间保证；若业务要求严格实时约束，需要额外的资源预留和可证明的最坏耗时控制。

\paragraph{不可行任务的处理}
当可信资源不可用、必需能力缺失或预算不足时，优先生成不可行原因，而不是在多个禁止动作中选取“分数最高者”。允许等待的任务设置最大等待窗口，等待时间计入约束成本；可以拆分的任务交还规划器重构；需授权或需求降级的任务暂停。任何新授权或降级都更新状态与约束版本，再重新求取可行域。

\subsection{策略网络、价值评估与约束学习}

\paragraph{状态编码与掩码策略}
状态编码器将任务语义向量与画像、资源状态、剩余预算和历史统计拼接。BGE类模型可提供语义向量，MLP输出候选动作分数$f_{\theta}(s,a)$。对安全可行域内的动作进行归一化：
\begin{equation}
\pi_{\theta}^{\mathrm{safe}}(a\mid s)=
\begin{cases}
\displaystyle\frac{\exp f_{\theta}(s,a)}
{\sum_{a'\in\mathcal{A}_{\mathrm{safe}}(s)}\exp f_{\theta}(s,a')},
&a\in\mathcal{A}_{\mathrm{safe}}(s),\\[7pt]
0,&\text{其他动作}.
\end{cases}
\label{eq:placement_masked_policy}
\end{equation}
可行域为空时不计算该分布，而进入不可行任务处理。掩码仅对已编码且使用最新数据检查的条件有效，不替代策略覆盖检查或发送前核验。

价值评估器分别估计收益价值$V_R(s)$与各成本价值$V_i(s)$。独立输出使费用、时间和失败预算具有可解释的预测，不需要将所有指标固定揉成单个静态权重分数。预测用于学习和决策，其准确性需要在独立轨迹上校准。

\paragraph{拉格朗日Actor--Critic更新}
对式~\eqref{eq:placement_cmdp}构造
\begin{equation}
\mathcal{L}(\theta,\boldsymbol{\lambda})=
J_R(\pi_\theta)-\sum_i\lambda_i
\left(J_i(\pi_\theta)-d_i\right),
\qquad\lambda_i\geq0.
\label{eq:placement_lagrangian}
\end{equation}
策略参数沿收益增加方向更新，约束乘子在观测成本超预算时上调。用收益与成本优势估计构造策略优势：
\begin{equation}
\widehat{A}^{\lambda}_t=
\widehat{A}_{R,t}-\sum_i\lambda_i\widehat{A}_{i,t}.
\end{equation}
更新形式为
\begin{equation}
\begin{aligned}
\theta&\leftarrow\theta+eta_{\theta}
\widehat{\nabla}_{\theta}\mathcal{L},\\
\lambda_i&\leftarrow
\left[\lambda_i+\eta_{\lambda}
\left(\widehat{J}_i-d_i\right)\right]_+.
\end{aligned}
\end{equation}
其中$\widehat{J}_i$应由完整回合或一致的截断引导估计得到，不能直接拿一个决策步的成本与整项任务预算比较。奖励与成本的尺度分别固定或归一化，乘子变化用于调整相应预算的约束压力。

训练时按实际掩码后的策略采样并记录动作概率，使用匹配行为策略的轨迹进行Actor--Critic更新。若调用前校验拒绝动作，将拒绝原因和时间开销记录为环境反馈；不将规则替换后的动作伪记为原策略采样结果。推理部署可使用确定性选择，但需独立评价该部署策略，因为训练期随机策略的期望约束不自动适用于取最大分数的策略。

拉格朗日方法可以出现振荡、估计误差与约束违反，需要通过独立验证和运行时硬过滤控制风险。这里提出的是可实现的学习与控制架构，不声称有限样本下收敛到全局最优，也不声称训练阶段始终满足全部累计预算。

\paragraph{训练数据与逐步接入}
首先由现有规则、分类路由和受控对照执行形成启动轨迹，记录任务条件、候选集、实际动作、质量证据、费用、时间和失败。监督预训练可让策略模仿合理路由或预测各候选收益，但这一阶段仍属于监督学习。

随后在沙箱任务或可控试验环境中，用掩码后的策略收集完整回合，依据式~\eqref{eq:placement_lagrangian}更新策略、价值网络和约束乘子。仅依赖旧策略日志训练时，必须处理动作覆盖和离策略偏差；本节采用先预训练、再受控在线采样的接入路径，不将直接套用在策略梯度上的旧日志当作有效强化学习训练。

训练、验证和测试按任务来源及时间区间划分；资源故障和网络扰动在测试中独立设置。调度器只读取当前可见信息，最终质量仅作为事后奖励，不能泄漏到当次状态。需要保存模型、画像规则、约束预算和策略版本，确保调度轨迹可复核。

\subsection{执行器分发、失败反馈与运行时接入}

选定动作被封装为SchedulingDecision和ResourceAllocation，包含具体端点、模型、策略版本、动作掩码、阶段计划、安全检查及选择依据。ExecutorRegistry据此调用本地、边缘HTTP或云侧兼容执行器，并统一返回成功状态、实际用量、耗时、错误类型和重试属性。实际推理位置以执行记录为准，不能仅用计划中写入的层级当作已发生的部署。

请求发出后，资源监控器更新端点级的健康度、限流与近期时延。限流端点进入受控冷却窗口；下一次调度读取新状态并再次过滤和选取。连接错误不一定意味着同层所有端点均失效，监控粒度应与动作中的具体端点对应。

\begin{figure}[htbp]
\centering
\fbox{\parbox{0.92\textwidth}{\centering
公开高复杂请求 $\rightarrow$ 选择云端点1 $\rightarrow$ 收到429\\[3pt]
$\rightarrow$ 计入本次费用、时间与失败成本\\[3pt]
$\rightarrow$ 更新端点1状态，重算剩余预算及可行域\\[3pt]
$\rightarrow$ 选择另一可行端点，或因能力/预算不足退出}}
\caption{将失败计入约束反馈后的有限重调度示例}
\label{fig:placement-reschedule}
\end{figure}

可重试错误在次数、剩余预算和截止时间允许时重新调度；不可重试错误直接返回处理原因。失败尝试也消耗预算，不能只统计最终成功调用。敏感任务的云端失败回退和边缘失败回退均需重新检查授权，恢复路径不能绕过原始约束。

运行时由HookManager在节点执行前获取共享状态与资源状态，调用调度接口并写入分配结果。节点结束后回传产物验收和实际消耗；任务最终质量可能延迟到质量门或结果汇聚后才获得，应通过任务和尝试标识归入对应轨迹，避免把当前HTTP成功直接当作任务成功奖励。

\begin{table}[htbp]
\centering
\caption{学习型调度轨迹的主要记录}
\label{tab:scheduling-trace}
\begin{tabular}{p{0.26\textwidth}p{0.63\textwidth}}
\toprule
类别 & 主要内容 \\
\midrule
决策输入 & 任务与资源画像、计划版本、剩余预算和可见历史 \\
可行域 & 候选动作、排除原因、可信域及能力检查结果 \\
策略输出 & 策略版本、动作概率或部署分数、价值预测和约束乘子 \\
实际执行 & 具体端点及模型、安全载荷校验、阶段计划和尝试标识 \\
执行反馈 & 质量证据、实际费用、墙钟时间、错误与失败成本 \\
预算与终止 & 剩余额度、超限事件、重试原因、不可行或暂停原因 \\
\bottomrule
\end{tabular}
\end{table}

\subsection{核心算法、实现边界与验证方案}

\begin{algorithm}[htbp]
\caption{硬可行域约束下的资源调度与执行}
\label{alg:tide-placement}
\KwIn{当前任务与共享状态$S_t$，资源监控$\Omega_t$，策略$\pi_\theta$}
\KwOut{执行结果、下一状态及决策轨迹}
编码任务、资源、历史与剩余预算，得到$s_t$\;
按授权、可信域、能力、可用性构造$\mathcal{A}_{\mathrm{safe}}(s_t)$\;
结合剩余预算排除不允许发起的调用\;
\If{无可行动作}{返回有界等待、拆分或暂停原因\;}
在实际可行域内按当前部署规则选择动作$a_t$\;
生成阶段计划，对实际发送载荷和最新资源状态再次校验\;
\If{校验失败}{记录原因与消耗，返回重新决策请求\;}
执行请求，记录实际用量、耗时和失败\;
更新资源监控、预算、任务状态和质量反馈关联\;
\Return 结果及$(s_t,a_t,r_t,\mathbf{c}_t,s_{t+1})$轨迹记录\;
\end{algorithm}

\begin{algorithm}[htbp]
\caption{累计预算约束下的有限重调度}
\label{alg:resilient-inference}
\KwIn{请求$q$，最大尝试次数$K$，剩余预算$b$}
\KwOut{结果或明确终止原因}
\For{$k\leftarrow1$ \KwTo $K$}{
    读取最新状态与约束，检查等待窗口及剩余预算\;
    调用算法~\ref{alg:tide-placement}并记录本次实际消耗\;
    \If{结果通过当前验收}{\Return 结果\;}
    \If{错误不可重试或预算不足或无可行资源}{\Return 终止原因\;}
    将失败写回端点状态，更新下一次调度输入\;
}
\Return 尝试预算耗尽及全部记录\;
\end{algorithm}

设候选动作数为$N$，任务编码成本为$T_{\mathrm{enc}}$，策略前向成本为$T_{\pi}$，资源检查及预算过滤为$O(N)$，单次决策成本为$O(T_{\mathrm{enc}}+T_{\pi}+N)$。最多$K$次重调度时，结构性决策开销相应乘以$K$；安全处理、质量评审、网络与模型推理耗时另计。训练成本由轨迹采集量、回合长度、模型规模和更新次数决定，不能用在线选择复杂度替代。

现有工程包含任务画像、规则与分类路由、可信域策略、执行器分发、资源监控和有限重调度。有限时域CMDP、掩码策略网络、成本价值网络、约束乘子训练及完整奖励归因为本节扩展设计。BGE+MLP可复用为网络结构，但需要完成策略训练与部署验证，才可报告约束强化学习的性能收益。

验证应在相同模型、任务与故障轨迹下，对比固定端点、规则调度、原有分类路由、无约束学习及完整约束学习。报告最终任务质量、费用、端到端时延、失败与期望预算违反，并另外报告单条轨迹超限率和硬约束违规。改变预算阈值形成质量--成本和质量--时延曲线，检查策略是否随预算调整，而不是只在一个固定权重下表现较好。

消融包括去掉剩余预算状态、约束乘子、硬动作过滤、失败反馈和质量奖励。硬约束消融仅在合成数据与隔离环境中验证行为。无可信候选、模型容量不足、限流、预算突然缩减和质量反馈延迟用于检验不可行退出及分布变化。所有扩展收益以实际完成的对照实验为依据，既有规则分支测试不作为学习收敛或最优性证据。

